\documentclass[12pt]{article}
\usepackage{fullpage}
\usepackage{amsmath,amssymb,amsthm}

\newcommand{\cL}{\mathcal L}
\newcommand{\Dr}{\operatorname{Dr}}
\newcommand{\Trop}{\operatorname{Trop}}
\newcommand{\rk}{\operatorname{rk}}
\newcommand{\cl}{\operatorname{cl}}
\newcommand{\op}{\operatorname{op}}

\newtheorem{theorem}{Theorem}
\newtheorem{lemma}{Lemma}
\newtheorem{proposition}{Proposition}

\begin{document}

\begin{center}
{\bf A canonical opposition involution on the relative Dressian}
\end{center}

\section*{Conventions}

Let $E$ be the ground set, let $r=\rk M$, and let $\mu_M$ be the
trivial rank $r$ valuated matroid supported on the bases of $M$.  I
use the standard projective-valuated-matroid model of the Dressian:
a rank $k$ point of $\Dr(M)$ is a rank $k$ valuated matroid
$\nu\preceq\mu_M$, modulo adding a scalar to all finite Plucker
coordinates, and
\[
        \Trop(\theta)\subseteq \Trop(\nu)
        \quad\Longleftrightarrow\quad
        \theta\preceq\nu .
\]
Here $p\preceq q$ means that $p$ is a valuated quotient of $q$.  For
valuated matroids $p,q$ of ranks $a\le b$ this quotient relation is
characterized by the nested tropical incidence relations
\[
 \min_{i\in T\setminus S}\{p(S\cup i)+q(T\setminus i)\}
 \quad\hbox{attained at least twice}                       \tag{I}
\]
for all $S\subseteq T$ with $|S|=a-1$ and $|T|=b+1$, together with the
valuated matroid axioms and the requirement that the underlying
ordinary matroid of $p$ is a quotient of the underlying ordinary matroid
of $q$.  This is a standard cryptomorphic form of the
incidence-Plucker criterion for flag valuated matroids; it is equivalent
to the arbitrary-index incidence relations often used in the literature.
If all displayed terms are $\infty$, the minimum is regarded as attained
at every index.

We also use the complete-flag cryptomorphism: a sequence
$g_0,\ldots,g_r$ of nonzero functions on the $j$-subsets, with the
trivial rank $0$ endpoint and top endpoint $\mu_M$, is a complete
valuated flag if and only if every adjacent pair satisfies (I).  This
is the homogeneous-layer form of Murota's $M^\natural$-concavity
criterion.

We shall use one standard completion fact.  If
$\nu\preceq\mu_M$ has rank $k$, then $\nu$ can be inserted into a
complete valuated flag ending at $\mu_M$.  Explicitly, up to independent
scalar normalizations of the layers, one may take
\[
 h_j(B)=
 \begin{cases}
 \min\{\nu(C): B\subseteq C,\ |C|=k\}, & j\le k,\ |B|=j,\\
 \min\{\nu(C): C\subseteq B,\ |C|=k\}, & j\ge k,\ B\in\mathcal I_j(M),\\
 \infty, & j\ge k,\ B\notin\mathcal I_j(M).
 \end{cases}                                                \tag{1}
\]
This is Murota's valuated truncation/elongation, equivalently the
homogeneous form of $M^\natural$-concavity.  The only endpoint
worth spelling out is $h_r=\mu_M$: if
$m(B)=\min\{\nu(C):C\subseteq B,\ |C|=k\}$ for an $M$-basis $B$, then
$m$ is constant on the basis graph of $M$.  Indeed, for adjacent bases
$B'=B-e+f$, choose $C\subseteq B$ minimizing $m(B)$.  If $e\notin C$
then $C\subseteq B'$.  If $e\in C$, apply (I) to
$\nu\preceq\mu_M$ with lower set $C-e$ and upper set $B\cup\{f\}$;
some finite term other than the $e$-term has value at most $\nu(C)$
and gives a $k$-subset of $B'$.  Thus $m(B')\le m(B)$, and symmetry
gives equality.

\section*{1. Flat-lattice consequences}

\begin{lemma}\label{lem:modular}
If $L=\cL(M)$ admits an order-reversing involution
$x\mapsto x^\perp$, then $L$ is modular and
$\rk(x^\perp)=r-\rk(x)$.
\end{lemma}

\begin{proof}
The flat lattice is finite, ranked, and semimodular.  An
anti-automorphism reverses maximal chains, so it sends rank $i$
elements to rank $r-i$ elements.  Applying it to the semimodular
inequality gives the reverse inequality; hence equality holds for
every pair of flats.  This is modularity.
\end{proof}

Thus $(x\vee y)^\perp=x^\perp\wedge y^\perp$ and
$(x\wedge y)^\perp=x^\perp\vee y^\perp$.

\begin{lemma}[flat constancy]\label{lem:flatconst}
Let $\nu\preceq \mu_M$ have rank $k$.  If $B,B'$ are independent
$k$-subsets of $M$ with $\cl(B)=\cl(B')$, then $\nu(B)=\nu(B')$.
Moreover $\nu(A)=\infty$ whenever $A$ is dependent in $M$.
\end{lemma}

\begin{proof}
The underlying ordinary matroid of $\nu$ is a quotient of $M$, so its
bases are independent in $M$; this gives the support assertion.  The
basis graph of the restriction of $M$ to a rank $k$ flat is connected.
It is therefore enough to compare $B=A\cup\{b\}$ and
$B'=A\cup\{b'\}$ with $|A|=k-1$ and
$\cl(Ab)=\cl(Ab')$.  Choose $C$ of size $r-k$ such that $AbC$ is an
$M$-basis; then $Ab'C$ is also an $M$-basis.  In (I) for
$\nu\preceq\mu_M$ with lower set $A$ and upper set
$A\cup\{b,b'\}\cup C$, the only finite $\mu_M$-terms are those indexed
by $b$ and $b'$, because deleting an element of $C$ leaves the
dependent set $Abb'$.  Hence $\nu(Ab)=\nu(Ab')$.
\end{proof}

We may write
\[
        \bar\nu(X)=\nu(B),\qquad \rk X=k,\quad \cl(B)=X.
\]

\section*{2. The opposition transform}

For a rank $k$ quotient $\nu\preceq\mu_M$ define a rank $r-k$
coordinate vector by
\[
 \nu^{\perp_M}(A)=
 \begin{cases}
   \bar\nu\bigl((\cl A)^\perp\bigr),
      & A\text{ independent in }M,\ |A|=r-k,\\
   \infty,&\text{otherwise.}
 \end{cases}                                                 \tag{2}
\]
Adding a scalar to $\nu$ adds the same scalar to $\nu^{\perp_M}$, and
(2) is visibly involutive on flat coordinates.

\begin{lemma}[elementary opposition]\label{lem:elementary}
Let $p\preceq q\preceq\mu_M$, with $\rk p=d$ and $\rk q=d+1$.  Then
$q^{\perp_M}$ and $p^{\perp_M}$ satisfy the adjacent incidence
relations of ranks $r-d-1$ and $r-d$.
\end{lemma}

\begin{proof}
Fix $A$ of size $r-d-2$ and distinct $x,y,z\notin A$; if $d=r-1$
there is nothing to check.  Put $U=\cl(A)$, $W=\cl(Axyz)$, and
$\delta=\rk W-\rk U$.  If $A$ is dependent, or if $\delta\le1$, all
three transformed terms are $\infty$.

Assume $\delta=2$.  Let $G=W^\perp$ and $H=U^\perp$, of ranks $d$ and
$d+2$.  Every term not forced to be $\infty$ has the common
$p$-coordinate $\bar p(G)$ and a $q$-coordinate
$\bar q(L_e)$ with
\[
        L_e=(\cl(Ae))^\perp,\qquad e\in\{x,y,z\}.
\]
In the rank-two interval $[G,H]$, either two of the non-forced terms
come from the same flat $L_e$, in which case they are equal and all
other terms are $\infty$, or the three relevant flats are distinct
points of $[G,H]$.  In the latter case choose a basis $S$ of $G$,
representatives $\ell_x,\ell_y,\ell_z$ in the three points, and a
complement $C$ extending $H$ to a basis of $M$.  The relation (I) for
$q\preceq\mu_M$ with lower set $S$ and upper set
$S\cup\{\ell_x,\ell_y,\ell_z\}\cup C$ has exactly the three finite
$\mu_M$-terms $\bar q(L_x),\bar q(L_y),\bar q(L_z)$.  Their minimum is
attained at least twice; adding $\bar p(G)$ proves the desired
three-term relation.

Assume $\delta=3$.  Then $G=W^\perp$ and $H=U^\perp$ have ranks
$d-1$ and $d+2$.  Set
\[
 P_x=(\cl(Ayz))^\perp,\quad P_y=(\cl(Axz))^\perp,\quad
 P_z=(\cl(Axy))^\perp
\]
and
\[
 Q_x=(\cl(Ax))^\perp,\quad Q_y=(\cl(Ay))^\perp,\quad
 Q_z=(\cl(Az))^\perp .
\]
The elements $x,y,z$ form a rank-three frame over $U$; by modularity
and anti-isomorphism, $P_x,P_y,P_z$ are the three atoms over $G$ in
the dual frame and
\[
        Q_x=P_y\vee P_z,\quad Q_y=P_x\vee P_z,\quad
        Q_z=P_x\vee P_y .
\]
Choose a basis $S$ of $G$ and representatives
$a_x\in P_x\setminus G$, $a_y\in P_y\setminus G$,
$a_z\in P_z\setminus G$.  The adjacent relation (I) for
$p\preceq q$ applied to $S$ and $a_x,a_y,a_z$ is exactly the
transformed three-term relation.
\end{proof}

\begin{proposition}\label{prop:set}
The formula $\nu\mapsto\nu^{\perp_M}$ defines an involution of the set
underlying $\Dr(M)$.
\end{proposition}

\begin{proof}
Insert $\nu$ into a complete flag
$h_0\preceq h_1\preceq\cdots\preceq h_r=\mu_M$ using (1).  Applying
Lemma \ref{lem:elementary} to each adjacent pair shows that the
reversed sequence
\[
        h_r^{\perp_M},h_{r-1}^{\perp_M},\ldots,h_0^{\perp_M}
\]
satisfies all adjacent incidence relations.  Its endpoints are the
rank $0$ trivial valuation and $\mu_M$.  By the complete-flag
criterion, it is a complete valuated flag.  Hence every layer
$h_j^{\perp_M}$, in particular $\nu^{\perp_M}$, is a quotient of
$\mu_M$.  Applying (2) twice recovers the original flat coordinates,
so the map is an involution, projectively well-defined.
\end{proof}

\begin{lemma}[ordinary polar quotients]\label{lem:ordinary}
Let $N$ be the underlying ordinary matroid of a quotient of $M$, of
rank $k$, and let $N^{\perp_M}$ be the ordinary support obtained from
(2).  On flats $X$ of $M$ its rank function is
\[
        \rho_{N^{\perp_M}}(X)=\rk_M(X)-k+\rho_N(X^\perp).       \tag{3}
\]
Consequently, if the underlying matroid of $p$ is a quotient of the
underlying matroid of $q$, then the underlying matroid of
$q^{\perp_M}$ is a quotient of the underlying matroid of
$p^{\perp_M}$.
\end{lemma}

\begin{proof}
Formula (3) is the usual polar-dual rank formula for matroids on a
modular geometric lattice; its bases are exactly the independent
$(r-k)$-sets whose polar flat spans $N$, which is the support of (2).
For the consequence, let $N_p,N_q$ have ranks $a,b$.  If $X\subseteq Y$
are flats of $M$, then
\[
\begin{aligned}
\rho_{N_q^{\perp_M}}(Y)-\rho_{N_q^{\perp_M}}(X)
 &= \rk(Y)-\rk(X)
    -\bigl(\rho_{N_q}(X^\perp)-\rho_{N_q}(Y^\perp)\bigr)\\
 &\le
    \rk(Y)-\rk(X)
    -\bigl(\rho_{N_p}(X^\perp)-\rho_{N_p}(Y^\perp)\bigr)\\
 &=\rho_{N_p^{\perp_M}}(Y)-\rho_{N_p^{\perp_M}}(X),
\end{aligned}
\]
because $N_p$ is a quotient of $N_q$.  This rank-increment inequality
is the ordinary strong-map criterion.
\end{proof}

\section*{3. Full order reversal}

We need one small consequence of being a quotient of $\mu_M$.

\begin{lemma}[circuit minima]\label{lem:circuit}
Let $w\preceq\mu_M$ have rank $k$, let $A$ be an independent
$(k-1)$-set of $M$, and let $C$ be a circuit of the contraction
$M/A$ with $|C|\ge2$.  Then
\[
        \min_{i\in C} w(A\cup i)
\]
is attained at least twice.
\end{lemma}

\begin{proof}
For each $i\in C$, the set $A\cup(C-i)$ is independent and spans the
same flat.  Choose a set $D$ such that $A\cup(C-i)\cup D$ is an
$M$-basis for one, hence every, $i\in C$.  Apply (I) to
$w\preceq\mu_M$ with lower set $A$ and upper set $A\cup C\cup D$.
The finite $\mu_M$-terms are precisely the terms indexed by
$i\in C$; deleting an element of $D$ leaves the circuit $C$ over $A$.
Thus the displayed minimum is attained at least twice.
\end{proof}

\begin{theorem}\label{thm:order}
If $p\preceq q\preceq\mu_M$, then
\[
        q^{\perp_M}\preceq p^{\perp_M}.
\]
Consequently (2) induces an order-reversing involution of $\Dr(M)$.
\end{theorem}

\begin{proof}
Let $\rk p=a$ and $\rk q=b$.  By Proposition \ref{prop:set}, both
transforms are valuated matroid quotients of $\mu_M$, and by
Lemma \ref{lem:ordinary} their underlying ordinary matroids are in the
required quotient order.  Thus, in the nested quotient criterion of the
conventions, it remains to check the incidence relations (I) for ranks
$r-b\le r-a$.  If $b=r$, the lower transformed rank is $0$ and there
are no relations.  Assume $b<r$.

Fix $A\subseteq T$ with
\[
        |A|=r-b-1,\qquad |T|=r-a+1,
\]
and set $B=T\setminus A$, $m=|B|=b-a+2$.  If $A$ is dependent, all
terms are $\infty$.  Otherwise put $U=\cl(A)$ and $W=\cl(T)$, so
$\rk U=r-b-1$.  A term can be finite only if $T-i$ is independent;
therefore, if $\delta=\rk W-\rk U\le m-2$, all terms are $\infty$.

First suppose $\delta=m$.  Then $T$ is independent (and this case is
empty when $a=0$).  Let $G=W^\perp$ and $H=U^\perp$, of ranks
$a-1$ and $b+1$.  For $i\in B$ define
\[
        X_i=(\cl(T-i))^\perp,\qquad
        Y_i=(\cl(Ai))^\perp .
\]
For $J\subseteq B$, the flats $\cl(A\cup J)$ have ranks
$\rk U+|J|$ and form the Boolean sublattice generated by this frame in
$[U,W]$.  After applying $^\perp$, the $X_i$ are the atoms over $G$ in
the dual frame and
\[
        Y_i=\bigvee_{j\in B,\ j\ne i} X_j .
\]
Choose a basis $S$ of $G$ and elements $x_i\in X_i\setminus G$.
Then $S\cup\{x_i:i\in B\}$ has size $b+1$, and the relation (I) for
$p\preceq q$ applied to this lower set $S$ and upper set
$S\cup\{x_i:i\in B\}$ has terms
\[
        p(Sx_i)+q\bigl(S\cup\{x_j:j\ne i\}\bigr)
        =\bar p(X_i)+\bar q(Y_i).
\]
These are exactly
\[
        p^{\perp_M}(T-i)+q^{\perp_M}(Ai),
\]
so the transformed incidence minimum is attained at least twice.

It remains to consider $\delta=m-1$.  In the contraction $M/A$, the
set $B$ has nullity one.  Let $C$ be its unique circuit.  If $C$ is a
singleton loop, then every term is $\infty$: the loop index makes
$Ai$ dependent, and every other deletion leaves the loop in $T-i$.
Thus assume $|C|\ge2$.  For $i\notin C$, the set $T-i$ is dependent,
so the transformed term is $\infty$.  For $i\in C$, the set $T-i$ is
independent with closure $W$, and hence
\[
        p^{\perp_M}(T-i)=\bar p(W^\perp)
\]
is independent of $i$.  The remaining factors are
$q^{\perp_M}(Ai)$.  Since $q^{\perp_M}\preceq\mu_M$ by
Proposition \ref{prop:set}, Lemma \ref{lem:circuit} applied to
$w=q^{\perp_M}$ and to the circuit $C$ of $M/A$ shows that
$\min_{i\in C}q^{\perp_M}(Ai)$ is attained at least twice.  Adding
the common summand $\bar p(W^\perp)$ proves the incidence relation.
\end{proof}

\section*{4. Agreement with the flat involution}

Let $F$ be a flat of rank $f$, and let $q_F$ be the trivial rank
$r-f$ valuated matroid of $M/F\oplus U_{0,F}$.  For an independent
$(r-f)$-set $B$ with $Z=\cl(B)$,
\[
        q_F(B)<\infty\quad\Longleftrightarrow\quad Z\vee F=\hat1.
                                                               \tag{4}
\]
For an independent $f$-set $A$ with $Y=\cl(A)$, (2) and (4) give
\[
 q_F^{\perp_M}(A)<\infty
 \quad\Longleftrightarrow\quad
        Y^\perp\vee F=\hat 1 .
\]
Since $\rk(Y^\perp)=r-f$ and $\rk F=f$, modularity makes this
equivalent to $Y^\perp\wedge F=\hat0$.  Applying the
order-reversing involution gives
\[
        Y\vee F^\perp=\hat1,
\]
which is precisely the basis condition for
$M/F^\perp\oplus U_{0,F^\perp}$.  The finite values are all zero on
both sides.  Therefore
\[
 \Phi\bigl(\Trop(M/F\oplus U_{0,F})\bigr)
 =
 \Trop(M/F^\perp\oplus U_{0,F^\perp}).
\]

Combining this with Theorem \ref{thm:order}, the answer to the
question is yes: the map represented by
\[
        \boxed{\ \nu\longmapsto \nu^{\perp_M}\ }
\]
is an order-reversing involution of $\Dr(M)$ extending
$F\mapsto F^\perp$.

\begin{thebibliography}{9}

\bibitem{BEZ}
M.~Brandt, C.~Eur and L.~Zhang,
\emph{Tropical flag varieties},
Advances in Mathematics \textbf{384} (2021), 107695.

\bibitem{DW}
A.~Dress and W.~Wenzel,
\emph{Valuated matroids},
Advances in Mathematics \textbf{93} (1992), 214--250.

\bibitem{JL}
M.~Jarra and O.~Lorscheid,
\emph{Flag matroids with coefficients},
Advances in Mathematics \textbf{436} (2024), 109396.

\bibitem{Hirai}
H.~Hirai,
\emph{Uniform semimodular lattices and valuated matroids},
Journal of Combinatorial Theory, Series A \textbf{165} (2019), 325--359.

\bibitem{Murota}
K.~Murota,
\emph{Matroid valuation on independent sets},
Journal of Combinatorial Theory, Series B \textbf{69} (1997), 59--78.

\bibitem{Speyer}
D.~Speyer,
\emph{Tropical linear spaces},
SIAM Journal on Discrete Mathematics \textbf{22} (2008), 1527--1558.

\end{thebibliography}

\end{document}
